
\documentclass[preprint,12pt]{interact}

\usepackage{amssymb}

\usepackage{amsmath}
\usepackage{bm}

\usepackage{float}
\floatstyle{plaintop}
\restylefloat{table}
\usepackage{caption}
\captionsetup[table]{singlelinecheck=off, justification=raggedright}
\usepackage{tablefootnote}
\usepackage{threeparttable}
\usepackage{adjustbox}
\usepackage{booktabs}
\usepackage{multirow}
\usepackage{array}
\usepackage{subcaption}
\captionsetup[subfigure]{justification=centering}


\begin{document}

\begin{frontmatter}



\title{Ray tracing – Curved quadratic elements}



\end{frontmatter}

\section{Intersection problem}\label{intersect}

Surface position on the tetrahedron's face (weighted sum of shape functions):
\begin{align}
	P(g,h) = \sum_{i-0}^{5} N_i(g,h) \cdot f_{nodes_i}
\end{align}

Ray:
\begin{align}
	R(t) = \bm{o} + t \bm{d}
\end{align}
where $\bm{d}$ is the ray's direction, and $\bm{o}$ is the ray's origin

Face-ray intersection (3 unknowns):
\begin{align}
	\bm{o} + t \bm{d} - P(g,h) = 0
\end{align}

Nonlinear system of equations:
\begin{align}
	F(t,g,h) = \bm{o} + t \bm{d} - \sum_{i-0}^{5} N_i(g,h) \cdot f_{nodes_i} = \bm{0}
\end{align}

\section{Möller-Trumbore algorithm for linear triangle intersection (initial guess)}\label{MT}

\begin{align}
	\bm{E_1} = V_1 - V_0, \; \bm{E_2} = V_2 - V_0 \\
	\bm{p} = \bm{d} \times \bm{E_2} \\
	det = \bm{E_1} \cdot \bm{p} \\
	\bm{t_{vec}} = \bm{o} - V_0
\end{align}
Where $V_0$, $V_1$, $V_2$ are triangle vertices.

\begin{align}
	u = \frac{\bm{t_{vec}} \cdot \bm{p} }{det}
\end{align}

\begin{align}
	\bm{q} = \bm{t_{vec}} \times \bm{E_1}
\end{align}

\begin{align}
	v = \frac{\bm{d} \cdot \bm{q} }{det}
\end{align}

\begin{align}
	t = \frac{\bm{E_2} \cdot \bm{q} }{det}
\end{align}


\section{Newton-Raphson}\label{NR}

Jacobian of $F$:
\begin{align}
	\bm{J_F} = \left[  \frac{\partial F}{\partial f}, \frac{\partial F}{\partial g}, \frac{\partial F}{\partial h} \right] \\
	\bm{J_F} = \left[  \bm{d}, -\frac{\partial P}{\partial g}, -\frac{\partial P}{\partial h} \right]
\end{align}

\begin{align}
	\frac{\partial P}{\partial g} = \sum_{i-0}^{5} \frac{\partial N_i}{\partial g} \cdot f_{nodes_i} \\
	\frac{\partial P}{\partial h} = \sum_{i-0}^{5} \frac{\partial N_i}{\partial h} \cdot f_{nodes_i}
\end{align}

The update by solving a linear system:
\begin{align}
	\bm{\Delta} = \left[ \Delta t, \Delta g, \Delta h \right]^T \\
	\bm{J_F} \cdot \bm{\Delta} = -F(t, g, h)
\end{align}

\begin{align}
\begin{bmatrix}
	\bm{d}_x & -\frac{\partial P_x}{\partial g} & -\frac{\partial P_x}{\partial h} \\
	\bm{d}_y & -\frac{\partial P_y}{\partial g} & -\frac{\partial P_y}{\partial h} \\
	\bm{d}_z & -\frac{\partial P_z}{\partial g} & -\frac{\partial P_z}{\partial h} 
\end{bmatrix} 
\begin{bmatrix}
	\Delta t \\
	\Delta g \\
	\Delta h
\end{bmatrix} = - (\bm{o} + t \bm{d} - P(g,h))
\end{align}



\end{document}

\endinput

